\usepackage[utf8]{inputenc}
\usepackage[german]{babel} % Deutsche Sprache 
\usepackage{amsmath}
\usepackage{graphicx}
\usepackage{fancyhdr}   % Kopf- / Fußzeile
\usepackage{geometry}
\usepackage{tikz}
\usepackage{siunitx}    % SI-Einheiten 
\usepackage{multirow}   % Tabellen
\usepackage{booktabs}   % Tabellen
\usepackage{siunitx}    % große Tabellen
\usepackage{hyperref}
\usepackage{listings}   % Code
\usepackage[framemethod=default]{mdframed} % Custom boxen
\usepackage{xcolor}
\usepackage{subcaption} % Mehrere Bilder nebeneinander
\usepackage[backend=biber,style=numeric, sorting=none]{biblatex} % Literaturverzeichnis
\usepackage{float}
\addbibresource{quellen.bib}

\usepackage[nameinlink]{cleveref} % Referrenzen Test
\renewcommand{\lstlistingname}{Code}             % für lstlistings umbennen zu Code
\crefname{listing}{Code}{Codes}
\Crefname{listing}{Code}{Codes} % für Satzanfang




\lstdefinelanguage{Kotlin}{
  morekeywords={val,var,fun,for,if,else,return,null,true,false,class,object},
  sensitive=true,
  morecomment=[l]{//},
  morecomment=[s]{/*}{*/},
  morestring=[b]",
}

\lstset{
  language=Kotlin,
  basicstyle=\ttfamily\footnotesize,
  keywordstyle=\color{blue},
  commentstyle=\color{gray},
  stringstyle=\color{orange},
  breaklines=true,
  captionpos=b
}

\lstset{language=Python,
        frame=single,
        basicstyle=\ttfamily\footnotesize, % Noch kleinere Schrift
        keywordstyle=\color{blue},
        commentstyle=\color[rgb]{0,0.5,0},
        stringstyle=\color{orange},
        numbers=left,
        numberstyle=\tiny\color{gray},
        stepnumber=1,
        numbersep=10pt,
        tabsize=4,
        showspaces=false,
        showstringspaces=false,
        breaklines=true,
        morekeywords={self, None, sudo, apt, install},
        captionpos=b
        }



\definecolor{codegray}{gray}{0.9}
\newcommand{\code}[1]{\colorbox{codegray}{\texttt{#1}}}


\geometry{  
            a4paper, 
            left=25mm, 
            right=25mm, 
            top=20mm, 
            bottom=20mm, 
            headsep=5mm % Abstand Headerlinie <-> Sektion
            }
\usepackage{mathptmx} % Use Times New Roman font

\hypersetup{pdftitle={Bericht Niklas Felhauer},pdfauthor={Niklas Felhauer}}

% Header and Footer
\pagestyle{fancy}
\fancyhf{}
\fancyhead[L]{\nouppercase{\leftmark}}
\fancyhead[R]{\thepage}

% Schriftart
\usepackage{avant} 
\newcommand*{\avantstyle}{\fontfamily{pag}\selectfont}
\DeclareTextFontCommand{\avantfont}{\avantstyle}
\usepackage{plex-sans}
\renewcommand*\familydefault{\sfdefault}

% Exercise Box
% Custom Box (Exercise)
\newmdenv[skipabove=7pt,
skipbelow=7pt,
rightline=false,
leftline=true,
topline=false,
bottomline=false,
backgroundcolor=white,
linecolor=ThemeColor,
innerleftmargin=5pt,
innerrightmargin=5pt,
innertopmargin=5pt,
innerbottommargin=5pt,
leftmargin=0cm,
rightmargin=0cm,
linewidth=4pt]{eBox}

% Define a custom environment to include title
\newenvironment{exercise}[1]{%
  \begin{eBox}%
  \textbf{\textcolor{ThemeColor}{\large #1}}\\
}{%
  \end{eBox}%
}